\documentclass[12pt]{article}

\usepackage[a4paper, margin=1in]{geometry}
\usepackage{amsmath}
\usepackage{amssymb}

\title{Black-Scholes Option Pricing Model in Excel/VBA\\
Pedagogical Implementation with Full Greeks Computation}
\author{Adam El Gbouri}
\date{February 2026}

\begin{document}

\maketitle

\section{Project Overview}

This project consists of the development of a pedagogical option pricing tool implemented in Excel and automated using VBA. The objective is to compute the price of European call and put options under the Black-Scholes-Merton framework, together with their associated sensitivity measures (Greeks).

The file is designed primarily as an educational support tool to facilitate understanding of option pricing mechanics, continuous-time modeling, and risk sensitivities. It separates numerical computation from financial interpretation to enhance clarity.

When the user clicks on the \textit{Start} button, a VBA macro executes all pricing calculations automatically. Interpretations are displayed through Excel formulas, while explanatory roles are provided as static descriptive text.

\section{Mathematical Framework}

Under the Black-Scholes-Merton assumptions, the underlying asset follows the stochastic differential equation:

\[
dS_t = (r - q) S_t dt + \sigma S_t dW_t
\]

where:

\begin{itemize}
    \item $r$ is the risk-free interest rate,
    \item $q$ is the continuous dividend yield,
    \item $\sigma$ is the volatility,
    \item $W_t$ is a standard Brownian motion.
\end{itemize}

The European call and put prices are given by:

\[
C = S_0 e^{-qT} N(d_1) - K e^{-rT} N(d_2)
\]

\[
P = K e^{-rT} N(-d_2) - S_0 e^{-qT} N(-d_1)
\]

where:

\[
d_1 =
\frac{\ln(S_0/K) + (r - q + \frac{1}{2}\sigma^2)T}
{\sigma \sqrt{T}}
\]

\[
d_2 = d_1 - \sigma \sqrt{T}
\]

and $N(\cdot)$ denotes the cumulative distribution function of the standard normal distribution.

\section{Structure of the Excel Model}

The Excel file is organized into four main sections:

\begin{enumerate}
    \item Inputs
    \item Calculations
    \item Option Pricing
    \item Greeks
\end{enumerate}

\subsection*{Inputs}

The user provides:

\begin{itemize}
    \item Spot price $S$
    \item Strike price $K$
    \item Maturity $T$
    \item Risk-free rate $r$
    \item Volatility $\sigma$
    \item Dividend yield $q$
\end{itemize}

\subsection*{Automated Computation}

When the \textit{Start} button is clicked, the VBA macro:

\begin{itemize}
    \item Computes $d_1$ and $d_2$
    \item Evaluates normal cumulative probabilities
    \item Computes call and put prices
    \item Calculates all primary Greeks
    \item Writes results automatically into designated cells
\end{itemize}

Interpretations are generated through Excel formulas to maintain transparency between computation and financial meaning.

\section{Greeks Computation}

The following sensitivities are computed:

\subsection*{Delta}

\[
\Delta_{call} = e^{-qT} N(d_1)
\]

\[
\Delta_{put} = \Delta_{call} - e^{-qT}
\]

Delta measures the sensitivity of the option price to changes in the underlying price.

\subsection*{Gamma}

\[
\Gamma =
\frac{e^{-qT} \phi(d_1)}
{S_0 \sigma \sqrt{T}}
\]

Gamma measures the curvature of the option price with respect to the underlying price.

\subsection*{Vega}

\[
Vega =
S_0 e^{-qT} \sqrt{T} \phi(d_1)
\]

Vega measures sensitivity to volatility.

\subsection*{Theta}

Theta measures time decay and is computed separately for calls and puts.

\subsection*{Rho}

\[
\rho_{call} =
K T e^{-rT} N(d_2)
\]

\[
\rho_{put} =
- K T e^{-rT} N(-d_2)
\]

Rho measures sensitivity to interest rates.

\section{Pedagogical Orientation}

The file is intentionally designed as an educational framework rather than a production trading tool. The objectives are:

\begin{itemize}
    \item To clarify the role of $d_1$ and $d_2$
    \item To illustrate the relationship between $N(d_1)$ and Delta
    \item To explain the economic meaning of each Greek
    \item To provide a clear separation between numerical computation and financial interpretation
\end{itemize}

The visual layout emphasizes understanding rather than computational complexity.

\section{Potential Enhancements}

Several extensions could transform this pedagogical tool into a more advanced financial application:

\subsection*{1. Interactive Input Interface}

Instead of manual cell inputs, a VBA UserForm (input dialog box) could be implemented. Upon clicking \textit{Start}, a graphical input window could allow the user to enter:

\begin{itemize}
    \item Asset symbol
    \item Strike
    \item Maturity
\end{itemize}

\subsection*{2. Market Data Integration}

Spot price, volatility, dividend yield, and risk-free rate could be automatically retrieved from external data providers such as:

\begin{itemize}
    \item Bloomberg API
    \item Refinitiv
    \item Yahoo Finance (via VBA web requests)
\end{itemize}

This would transform the file into a semi-professional pricing tool.

\subsection*{3. Volatility Surface Integration}

Implied volatility could be extracted dynamically depending on strike and maturity.

\subsection*{4. Scenario and Sensitivity Analysis}

Data tables could be added to:

\begin{itemize}
    \item Analyze sensitivity to spot changes
    \item Study volatility shifts
    \item Evaluate time decay
\end{itemize}

\subsection*{5. Delta-Hedging Simulation}

A dynamic hedging module could simulate portfolio rebalancing and P\&L evolution.

\section{Conclusion}

This project provides a structured implementation of the Black-Scholes model within Excel using VBA automation. It combines theoretical rigor with practical usability and pedagogical clarity.

While the current version focuses on transparency and educational value, the architecture allows for further extensions toward real-time data integration and dynamic hedging analysis.

The framework demonstrates a solid understanding of continuous-time finance, along with stochastic modeling and risk sensitivity analysis.

\end{document}